
\maketitle

\begin{abstract}
We prove that a compatible family of finitely additive charges on a directed system
of Boolean algebras extends uniquely to a $\sigma$-additive measure on the generated
$\sigma$-algebra, using Stone duality as the principal mechanism.  The argument
proceeds in two steps.  First, Stone duality converts the directed system of Boolean
algebras into a cofiltered inverse system of compact Hausdorff spaces; compactness of
the inverse limit then yields a $\sigma$-additive measure on the Borel algebra of the
limit space without any countable additivity assumption on the input charges.  Second,
the Borel algebra of the inverse limit is identified with the observable $\sigma$-algebra
generated by the cylinder sets, under two conditions drawn from the broader
programme: discriminability of the query system and collective exhaustion of the
charge family.

The result gives a second, independent proof of the observable extension theorem
established in \cite{mahon_paper0} by Carath\'{e}odory methods.  The two routes
complement each other: the Carath\'{e}odory route assumes $\sigma$-additivity of the
marginals and derives the global measure directly; the Stone route assumes only finite
additivity and derives $\sigma$-additivity from topological compactness.  Under shared
hypotheses the two measures coincide.

The proof makes explicit a role for discriminability --- the condition that queries
separate points of the sample space --- as a structural prerequisite for the
identification of the observable $\sigma$-algebra with the Borel algebra of the Stone
space.  Collective exhaustion, previously identified as irreducible in
\cite{mahon_paper0}, reappears here as the condition that the Stone measure is
supported on the principal ultrafilters, i.e., on the image of the sample space
inside its Stone compactification.
\end{abstract}

% ------------------------------------------------------------------
\section{Introduction}
% ------------------------------------------------------------------

A central problem in the foundations of probability is the extension of consistent
finite-dimensional distributions to a global measure on an infinite-dimensional space.
The classical solution, due to Kolmogorov, requires the marginals to be $\sigma$-additive
from the outset and imposes topological conditions (standard Borel spaces) to guarantee
tightness.  A more recent approach, developed in \cite{mahon_paper0}, replaces these
topological hypotheses with structural conditions on a system of observable queries and
derives the global measure by Carath\'{e}odory extension.

The present paper offers a third route, using Stone duality.  The key insight is that
a Boolean algebra $B$ and its Stone space $\St(B)$ carry equivalent information: a
finitely additive charge on $B$ corresponds bijectively to a regular Borel measure on
the compact Hausdorff space $\St(B)$.  On compact spaces, regular Borel measures are
automatically $\sigma$-additive.  So finite additivity on a Boolean algebra, combined
with Stone duality, yields $\sigma$-additivity for free on the Stone space side.

When the Boolean algebras form a directed system, Stone duality converts this into an
inverse system of compact Hausdorff spaces.  The Stone space of the direct limit is the
inverse limit of the individual Stone spaces (Proposition~\ref{prop:stepA}).  A
compatible family of charges on the directed system therefore assembles into a
$\sigma$-additive measure on the inverse limit (Proposition~\ref{prop:invlim}).

The remaining task is to identify the Borel algebra of the inverse limit with the
observable $\sigma$-algebra $\sigma(\CylGen)$ generated by all cylinder sets.  This
identification has two parts.  The structural part (Step~B1,
Proposition~\ref{prop:B1}) uses the fact that the clopens of the Stone space pull back
under the Stone embedding $\pure : \Omega \to \St(\CylGen)$ to exactly the cylinder
sets; discriminability ensures the embedding is injective.  The measure-theoretic part
(Step~B2, Proposition~\ref{prop:B2}) uses collective exhaustion to show that no mass
escapes to the non-principal ultrafilters in $\St(\CylGen)$.

\paragraph{Relation to prior work.}
The individual components of this argument are known.  The correspondence between
charges on a Boolean algebra and measures on its Stone space is made explicit in
\cite{cardona2025}, Section~6.  The extension of compatible measures on an inverse
system of compact spaces is due to Choksi \cite{choksi1958}.  The Yosida--Hewitt
decomposition \cite{yosida1952} identifies the purely finitely additive component of a
charge with mass on non-principal ultrafilters.  What appears to be new is the
assembly of these components into a single theorem for directed systems of Boolean
algebras arising from query systems, with the precise roles of discriminability and
collective exhaustion identified.

\paragraph{Organisation.}
Section~\ref{sec:setup} fixes notation and states the three hypotheses.
Section~\ref{sec:stepA} proves Step~A.  Section~\ref{sec:invlim} constructs the
inverse limit measure.  Sections~\ref{sec:B1} and~\ref{sec:B2} prove the two parts of
Step~B.  Section~\ref{sec:assembly} assembles the main theorem.

% ------------------------------------------------------------------
\section{Setup and hypotheses}
\label{sec:setup}
% ------------------------------------------------------------------

\subsection{Query systems}

\begin{definition}[Query system]
A \emph{query system} is a tuple
$(\iota, \leq, \Omega, \{(\mathrm{Out}(i), \Fcal_i)\}_{i \in \iota},
\{\mathrm{eval}_i\}_{i \in \iota}, \{\pi_{ij}\}_{i \leq j})$
consisting of:
\begin{itemize}[noitemsep]
  \item a nonempty partially ordered set $(\iota, \leq)$, the \emph{index set};
  \item a set $\Omega$, the \emph{sample space};
  \item for each $i \in \iota$, a measurable space $(\mathrm{Out}(i), \Fcal_i)$,
    the \emph{outcome space at level} $i$;
  \item for each $i \in \iota$, a map $\mathrm{eval}_i : \Omega \to \mathrm{Out}(i)$,
    the \emph{evaluation map at level} $i$;
  \item for each $i \leq j$ in $\iota$, a measurable surjection
    $\pi_{ij} : \mathrm{Out}(j) \to \mathrm{Out}(i)$, the \emph{refinement map},
\end{itemize}
satisfying the \emph{coherence condition}: $\pi_{ij} \circ \mathrm{eval}_j =
\mathrm{eval}_i$ for all $i \leq j$.
\end{definition}

\subsection{Cylinder sets and the observable \texorpdfstring{$\sigma$}{sigma}-algebra}

\begin{definition}[Cylinder sets]
For $i \in \iota$ and $A \in \Fcal_i$, the \emph{level-$i$ cylinder with base $A$} is
\[
  \Cyl{i}{A} \;=\; \{\,\omega \in \Omega : \mathrm{eval}_i(\omega) \in A\,\}.
\]
The \emph{cylinder generating family} is
\[
  \CylGen \;=\; \bigl\{\,\Cyl{i}{A} : i \in \iota,\; A \in \Fcal_i\,\bigr\}.
\]
For each $i \in \iota$, let $B_i$ denote the Boolean algebra of subsets of $\Omega$
generated by $\{\Cyl{i}{A} : A \in \Fcal_i\}$.  Then $\CylGen = \bigcup_{i \in \iota} B_i$.

The \emph{observable $\sigma$-algebra} is $\sigma(\CylGen)$, the $\sigma$-algebra on
$\Omega$ generated by all cylinder sets.
\end{definition}

\begin{remark}
The coherence condition gives, for $i \leq j$,
\[
  \Cyl{i}{A} \;=\; \Cyl{j}{\pi_{ij}^{-1}(A)},
\]
so the same subset of $\Omega$ can be described as a cylinder at any level $j \geq i$.
\end{remark}

\subsection{Connecting maps}

For $i \leq j$, define the \emph{connecting map}
\[
  \varphi_{ij} : B_i \to B_j, \qquad
  \varphi_{ij}(\Cyl{i}{A}) \;=\; \Cyl{j}{\pi_{ij}^{-1}(A)}.
\]
This is a well-defined injective Boolean algebra homomorphism (Lemma~\ref{lem:phi_inj}
below).  The family $\{B_i, \varphi_{ij}\}$ is a directed system of Boolean algebras
with injective connecting maps.

\subsection{Compatible charges}

\begin{definition}[Charge, compatible family]
A \emph{charge} on a Boolean algebra $B$ is a finitely additive function
$\mu : B \to [0,\infty)$ with $\mu(\emptyset) = 0$.

A family of charges $\{\mu_i\}_{i \in \iota}$ with $\mu_i$ on $B_i$ is
\emph{compatible} if
\[
  \mu_j(\varphi_{ij}(E)) \;=\; \mu_i(E)
  \qquad \text{for all } i \leq j \text{ and all } E \in B_i.
\]
\end{definition}

\subsection{The three hypotheses}

The main theorem requires three conditions on the query system and charge family.

\begin{definition}[EvalSurjective]
The query system satisfies \emph{EvalSurjective} if
$\mathrm{eval}_i : \Omega \to \mathrm{Out}(i)$ is surjective for every $i \in \iota$.
\end{definition}

\begin{definition}[Discriminability]
The query system \emph{discriminates points} if for every $\omega \neq \omega'$ in
$\Omega$ there exists $E \in \CylGen$ with $\omega \in E$ and $\omega' \notin E$.
\end{definition}

\begin{definition}[Collective Exhaustion]
A charge $\mu$ on $\CylGen$ (the common extension of a compatible family $\{\mu_i\}$)
satisfies \emph{collective exhaustion} (CE) if: whenever $E_n \in \CylGen$ and
$E_1 \supseteq E_2 \supseteq \cdots$ with $\bigcap_{n} E_n = \emptyset$, we have
$\mu(E_n) \to 0$.
\end{definition}

\begin{remark}
CE is not derivable from any structural or algebraic condition on the query system
alone \cite{mahon_paper0}.  It is an irreducible condition on the charge family: it
rules out purely finitely additive charges, which satisfy every finite consistency
condition while assigning persistent mass to sequences that vanish in the limit.
\end{remark}

% ------------------------------------------------------------------
\section{Step A: the Stone space of the direct limit}
\label{sec:stepA}
% ------------------------------------------------------------------

We first verify that the connecting maps are injective, then invoke the standard
Stone duality fact.

\begin{lemma}[Injectivity of connecting maps]
\label{lem:phi_inj}
If the query system satisfies \emph{EvalSurjective}, then $\varphi_{ij} : B_i \to B_j$
is an injective Boolean algebra homomorphism for every $i \leq j$.
\end{lemma}

\begin{proof}
\textbf{Step 1: $\pi_{ij}$ is surjective.}
Let $x \in \mathrm{Out}(i)$.  Since $\mathrm{eval}_i$ is surjective, there exists
$\omega \in \Omega$ with $\mathrm{eval}_i(\omega) = x$.  By coherence,
$\pi_{ij}(\mathrm{eval}_j(\omega)) = \mathrm{eval}_i(\omega) = x$, so $x \in
\mathrm{im}(\pi_{ij})$.

\textbf{Step 2: $\varphi_{ij}$ is injective.}
Suppose $\varphi_{ij}(\Cyl{i}{A}) = \varphi_{ij}(\Cyl{i}{A'})$, i.e.,
$\Cyl{j}{\pi_{ij}^{-1}(A)} = \Cyl{j}{\pi_{ij}^{-1}(A')}$.  This gives
$\pi_{ij}^{-1}(A) = \pi_{ij}^{-1}(A')$.  Since $\pi_{ij}$ is surjective, the
preimage map $\pi_{ij}^{-1}$ is injective on subsets of $\mathrm{Out}(i)$, so
$A = A'$, hence $\Cyl{i}{A} = \Cyl{i}{A'}$.  Injectivity on generators extends to
the Boolean algebra $B_i$ by the universal property.
\end{proof}

\begin{proposition}[Step A]
\label{prop:stepA}
The Stone space of the direct limit $\St\!\bigl(\bigcup_i B_i\bigr)$ is homeomorphic
to the inverse limit $\limproj_i \St(B_i)$.
\end{proposition}

\begin{proof}
By Lemma~\ref{lem:phi_inj}, the connecting maps $\varphi_{ij}$ are injective Boolean
algebra homomorphisms.  Under Stone duality, injective Boolean algebra homomorphisms
correspond to surjective continuous maps between Stone spaces.  A directed system of
Boolean algebras with injective connecting maps therefore dualises to a cofiltered
inverse system of compact Hausdorff spaces with surjective bonding maps.

The Stone space functor $\St$ converts colimits (direct limits) in the category of
Boolean algebras to limits (inverse limits) in the category of compact Hausdorff
spaces \cite[II.4]{johnstone1982}.  Applied to $\bigcup_i B_i$ with its injective
system, this gives $\St\!\bigl(\bigcup_i B_i\bigr) \cong \limproj_i \St(B_i)$.
\end{proof}

% ------------------------------------------------------------------
\section{The inverse limit measure}
\label{sec:invlim}
% ------------------------------------------------------------------

\begin{proposition}[Inverse limit measure]
\label{prop:invlim}
Let $\{\mu_i\}$ be a compatible family of normalised charges on $\{B_i\}$.  Then there
exists a unique probability measure $\hat{\mu}$ on $\Borel(\limproj_i \St(B_i))$ such
that the pushforward of $\hat{\mu}$ along each projection $\limproj_i \St(B_i) \to
\St(B_j)$ agrees with the Borel measure on $\St(B_j)$ corresponding to $\mu_j$.
\end{proposition}

\begin{proof}
\textbf{Step 1: Single-algebra correspondence.}
For each $i$, the Boolean algebra $B_i$ is isomorphic (via Stone duality) to the
clopen algebra $\Clop(\St(B_i))$ of its Stone space.  A charge $\mu_i$ on $B_i$
therefore defines a finitely additive set function on the clopens of the compact
Hausdorff space $\St(B_i)$.  Since $\St(B_i)$ is compact and totally disconnected,
the clopens form a base for the topology.  A finitely additive charge on a base of
clopens of a compact space extends uniquely to a regular Borel measure; see
\cite{cardona2025}, Section~6, Theorem~D (or \cite{halmos1950}, \S\S53--54 for
the content construction).  Call this measure $\hat{\mu}_i$ on $\St(B_i)$.

\textbf{Step 2: Compatibility.}
The compatibility condition $\mu_j(\varphi_{ij}(E)) = \mu_i(E)$ translates, under the
Stone correspondence, to the condition that $\hat{\mu}_i$ is the pushforward of
$\hat{\mu}_j$ along the bonding map $\St(B_j) \to \St(B_i)$.  Thus the family
$\{\hat{\mu}_i\}$ is a compatible family of Borel probability measures on the inverse
system $\{\St(B_i)\}$.

\textbf{Step 3: Extension to the inverse limit.}
Each $\St(B_i)$ is compact Hausdorff.  The inverse limit $\limproj_i \St(B_i)$ is
compact Hausdorff (being a closed subspace of a product of compact Hausdorff spaces).
By \cite[Theorem~1]{choksi1958}, a compatible family of regular Borel probability
measures on a cofiltered inverse system of compact Hausdorff spaces with surjective
bonding maps extends uniquely to a regular Borel probability measure on the inverse
limit.  This gives $\hat{\mu}$ on $\Borel(\limproj_i \St(B_i))$.
\end{proof}

% ------------------------------------------------------------------
\section{Step B1: the structural identification}
\label{sec:B1}
% ------------------------------------------------------------------

Let $\pure : \Omega \to \St(\CylGen)$ denote the Stone embedding, which sends
$\omega \in \Omega$ to its principal ultrafilter $\pure(\omega) = \{E \in \CylGen :
\omega \in E\}$.  For $E \in \CylGen$, write $[E] = \{u \in \St(\CylGen) : E \in u\}$
for the corresponding basic clopen in $\St(\CylGen)$.

\begin{proposition}[Step B1: structural identification]
\label{prop:B1}
Suppose the query system discriminates points.  Then
\[
  \pure^{-1}\!\bigl(\Borel(\St(\CylGen))\bigr) \;=\; \sigma(\CylGen).
\]
\end{proposition}

\begin{proof}
We show the two inclusions separately.

\textbf{($\supseteq$): every cylinder set is in the pullback.}
For any $E = \Cyl{i}{A} \in \CylGen$,
\[
  \pure^{-1}([E]) \;=\; \{\omega \in \Omega : E \in \pure(\omega)\}
  \;=\; \{\omega \in \Omega : \omega \in E\} \;=\; E.
\]
Since $[E]$ is a clopen in $\St(\CylGen)$, it is Borel.  Therefore every cylinder set
lies in $\pure^{-1}(\Borel(\St(\CylGen)))$, and since this pullback is a
$\sigma$-algebra, $\sigma(\CylGen) \subseteq \pure^{-1}(\Borel(\St(\CylGen)))$.

\textbf{($\subseteq$): every element of the pullback is in $\sigma(\CylGen)$.}
The Stone space $\St(\CylGen)$ is compact and totally disconnected, so its clopens
form a topological basis.  The clopens of $\St(\CylGen)$ are exactly the sets $[E]$
for $E \in \CylGen$ \cite[I.3.2]{johnstone1982}.  Since the topology is generated by
these clopens, the Baire $\sigma$-algebra of $\St(\CylGen)$ --- the smallest
$\sigma$-algebra making all continuous functions measurable --- coincides with the
$\sigma$-algebra generated by the clopens $\{[E] : E \in \CylGen\}$.

\begin{remark}
On a compact Hausdorff space the Baire $\sigma$-algebra is contained in the Borel
$\sigma$-algebra, with equality when the space is second-countable.  We work
throughout with the Baire $\sigma$-algebra on $\St(\CylGen)$; this suffices for the
measure-theoretic purposes of Section~\ref{sec:assembly} because the measure
$\hat{\mu}$ constructed in Section~\ref{sec:invlim} is a Baire measure (it is
the unique extension of a charge on the clopen algebra), and Choksi's theorem
\cite{choksi1958} applies to Baire measures on compact spaces.
\end{remark}

Now, for any Baire set $V \in \Borel_{\mathrm{Ba}}(\St(\CylGen))$, we have
$V \in \sigma(\{[E] : E \in \CylGen\})$, so $\pure^{-1}(V) \in
\sigma(\{\pure^{-1}([E]) : E \in \CylGen\}) = \sigma(\CylGen)$.

\textbf{Injectivity.}
Discriminability implies $\pure$ is injective: if $\omega \neq \omega'$, there exists
$E \in \CylGen$ with $\omega \in E$ and $\omega' \notin E$, so $E \in \pure(\omega)$
but $E \notin \pure(\omega')$, hence $\pure(\omega) \neq \pure(\omega')$.  Injectivity
ensures that the pullback does not conflate distinct events: $\pure^{-1}$ is
injective on sets in the range of $[\ \cdot\ ]$.
\end{proof}

% ------------------------------------------------------------------
\section{Step B2: CE and the support condition}
\label{sec:B2}
% ------------------------------------------------------------------

\begin{proposition}[Step B2: support condition]
\label{prop:B2}
Let $\mu$ be a charge on $\CylGen$ satisfying collective exhaustion.  Then the
corresponding Baire measure $\hat{\mu}$ on $\St(\CylGen)$ is supported on
$\pure(\Omega)$, the set of principal ultrafilters.
\end{proposition}

\begin{proof}
By the Yosida--Hewitt theorem \cite{yosida1952}, every bounded charge $\mu$ on a
Boolean algebra decomposes uniquely as
\[
  \mu \;=\; \mu_c + \mu_p,
\]
where $\mu_c$ is countably additive and $\mu_p$ is purely finitely additive (i.e.,
$0 \leq \nu \leq \mu_p$ and $\nu$ countably additive implies $\nu = 0$).

Under Stone duality, the countably additive part $\mu_c$ corresponds to mass
concentrated on the principal ultrafilters $\pure(\Omega) \subset \St(\CylGen)$, while
the purely finitely additive part $\mu_p$ corresponds to mass on the non-principal
ultrafilters --- the ``phantom'' points of the Stone compactification.  More precisely,
if $\hat{\mu} = \hat{\mu}_c + \hat{\mu}_p$ is the corresponding decomposition of the
Baire measure on $\St(\CylGen)$, then $\hat{\mu}_p$ is supported on
$\St(\CylGen) \setminus \pure(\Omega)$; see \cite{rao1983}, Chapter~10.

Collective exhaustion is equivalent to $\mu_p = 0$: CE requires that $\mu(E_n) \to 0$
whenever $E_n \searrow \emptyset$, which is exactly the condition that the purely
finitely additive part vanishes \cite{mahon_paper0}, Theorem~(sp1\_iff).  Therefore
$\hat{\mu}_p = 0$, and $\hat{\mu} = \hat{\mu}_c$ is supported on $\pure(\Omega)$.
\end{proof}

% ------------------------------------------------------------------
\section{The main theorem}
\label{sec:assembly}
% ------------------------------------------------------------------

\begin{theorem}[Stone extension theorem]
\label{thm:main}
Let $(\iota, \leq, \Omega, \{\mathrm{Out}(i)\}, \{\mathrm{eval}_i\}, \{\pi_{ij}\})$
be a query system satisfying \emph{EvalSurjective} and \emph{Discriminability}, and
let $\{\mu_i\}$ be a compatible family of normalised charges on $\{B_i\}$ satisfying
\emph{Collective Exhaustion}.  Then there exists a unique probability measure $P$ on
$(\Omega, \sigma(\CylGen))$ such that
\[
  P(\Cyl{i}{A}) \;=\; \mu_i(\Cyl{i}{A})
  \qquad \text{for all } i \in \iota \text{ and } A \in \Fcal_i.
\]
\end{theorem}

\begin{proof}
The proof assembles the preceding propositions along the following chain:
\[
  \{\mu_i\}
  \;\xrightarrow{\S\ref{sec:invlim},\,\text{Step 1}}\;
  \{\hat{\mu}_i\}
  \;\xrightarrow{\S\ref{sec:invlim},\,\text{Step 3}}\;
  \hat{\mu} \text{ on } \limproj_i \St(B_i)
  \;\xrightarrow{\text{B1}}\;
  P \text{ on } \sigma(\CylGen).
\]

\textbf{Existence.}
By Proposition~\ref{prop:stepA}, $\St(\CylGen) \cong \limproj_i \St(B_i)$.  By
Proposition~\ref{prop:invlim}, the compatible family $\{\mu_i\}$ extends to a Baire
probability measure $\hat{\mu}$ on $\limproj_i \St(B_i) \cong \St(\CylGen)$.

By Proposition~\ref{prop:B2}, $\hat{\mu}$ is supported on $\pure(\Omega)$.  Since
$\pure : \Omega \to \St(\CylGen)$ is injective (Discriminability), the pushforward
$P = (\pure)_* \hat{\mu}$ --- the image measure under $\pure^{-1}$ on the support ---
is a well-defined probability measure on $\Omega$.

By Proposition~\ref{prop:B1}, $\pure^{-1}$ maps Baire sets in $\St(\CylGen)$ to
sets in $\sigma(\CylGen)$, so $P$ is defined on $\sigma(\CylGen)$.  For any
cylinder set $E = \Cyl{i}{A}$,
\[
  P(E) \;=\; \hat{\mu}([E]) \;=\; \mu_i(\Cyl{i}{A}),
\]
where the last equality holds by the single-algebra Stone correspondence
(Proposition~\ref{prop:invlim}, Step~1).

\textbf{Uniqueness.}
Two probability measures on $(\Omega, \sigma(\CylGen))$ that agree on all cylinder
sets are equal, by the $\pi$-$\lambda$ theorem applied to the $\pi$-system $\CylGen$
(under the upper-directedness condition on $\iota$, which is implied by EvalSurjective
in the query system setting); see \cite{mahon_paper0}, Theorem~(observational\_determination).
\end{proof}

\begin{remark}[Relation to the Carath\'{e}odory route]
Under the additional assumption that each $\mu_i$ is already $\sigma$-additive (a
countably additive probability measure), the measure $P$ produced by
Theorem~\ref{thm:main} coincides with the measure produced by the Carath\'{e}odory
extension in \cite{mahon_paper0}.  Both measures agree on $\CylGen$; uniqueness
(observational determination) then forces them to be equal.  The Stone route thus
provides an independent proof of the observable extension theorem, requiring only
finite additivity of the input charges.
\end{remark}

\begin{remark}[CE is not bypassed]
It may appear that the compactness of the Stone space eliminates the need for
Collective Exhaustion.  Compactness does ensure that the measure $\hat{\mu}$ on
$\St(\CylGen)$ is $\sigma$-additive.  But $\hat{\mu}$ lives on the Stone space, not
on $\Omega$.  The descent from $\St(\CylGen)$ to $\Omega$ requires that no mass
escapes to the non-principal ultrafilters --- and that is precisely CE (Proposition
\ref{prop:B2}).  In the Stone picture, CE is not an additional hypothesis but a
support condition: it says the measure concentrates on the image of the sample space
inside its Stone compactification.
\end{remark}

\bibliographystyle{plain}
\bibliography{references}
